% Real vs estimated scale for the SAM 3 + VGGT-SLAM lists (internal). Fragment for % Real vs estimated scale for the SAM 3 + VGGT-SLAM lists (internal). Fragment for % Real vs estimated scale for the SAM 3 + VGGT-SLAM lists (internal). Fragment for % Real vs estimated scale for the SAM 3 + VGGT-SLAM lists (internal). Fragment for \input{scale}; needs \usepackage{booktabs}
% and \usepackage{url}. Numbers from 3d-scene-lang/data/scenelang_vsibench_slam/scale_compare.py (gandalf).

\section{Real vs.\ estimated scale}
\label{sec:scale}

We compare the SAM~3 + VGGT-SLAM box lists at real scale (ground-truth scale from the camera fit) and at estimated scale
(from class size priors). We use the SFT Qwen3.5-4B Exp1 model
(\path{sr_vggt_novggt_qwen35_4b_reason190k_mixv2_1_objq/checkpoint-569}) on all 5,130 VSI-Bench questions. Each scale is run
with two inference rules: mix-v2-1 (video only for route planning, room size and questions with no named object) and
the list on every question (``full''). Table~\ref{tab:scale} gives the scores per task, together with the share of answers that
change between real and estimated scale (full rule) and which side is better on the changed answers.
Answers are extracted as the last option letter or number after \texttt{\textless/think\textgreater}. This makes the absolute values differ
slightly from the main tables (e.g.\ 59.7 vs.\ 60.8 on absolute distance for real mix-v2-1); the differences between runs
are what matter here.

\begin{table}[h]
\centering
\caption{Exp1 model with SAM~3 + VGGT-SLAM lists at real vs.\ estimated scale, per task (MRA or accuracy, \%).
Changed: answers that differ between real and estimated scale (full rule); est./real better: of the changed answers, how many
score higher at each scale.}
\label{tab:scale}
\small
\begin{tabular}{lrrrrrrrr}
\toprule
 & & \multicolumn{2}{c}{mix-v2-1} & \multicolumn{2}{c}{full} & \multicolumn{3}{c}{real $\rightarrow$ est.\ (full)} \\
\cmidrule(lr){3-4}\cmidrule(lr){5-6}\cmidrule(lr){7-9}
Task & $n$ & real & est. & real & est. & changed & est.\ better & real better \\
\midrule
\multicolumn{9}{l}{\emph{Depends on scale}} \\
Absolute distance & 834 & 59.7 & 57.5 & 59.7 & 57.5 & 52.3\% & 163 & 193 \\
Object size       & 953 & 76.4 & 75.5 & 76.4 & 75.5 & 56.9\% & 153 & 142 \\
Room size         & 288 & 72.8 & 72.8 & 72.9 & 73.0 & 70.5\% & 78 & 68 \\
\midrule
\multicolumn{9}{l}{\emph{Independent of scale}} \\
Counting                   & 565 & 75.4 & 75.8 & 75.4 & 75.8 & 3.2\% & 9 & 6 \\
Relative distance          & 710 & 69.2 & 68.7 & 69.2 & 68.7 & 1.4\% & 3 & 6 \\
Relative direction, easy   & 217 & 78.8 & 78.3 & 78.8 & 78.3 & 1.4\% & 1 & 2 \\
Relative direction, medium & 378 & 79.9 & 79.6 & 79.9 & 79.6 & 0.8\% & 1 & 2 \\
Relative direction, hard   & 373 & 72.9 & 72.9 & 72.9 & 72.9 & 1.3\% & 2 & 2 \\
Appearance order           & 618 & 80.4 & 80.4 & 80.7 & 80.7 & 0.3\% & 1 & 1 \\
Route planning             & 194 & 49.5 & 49.5 & 47.9 & 49.5 & 4.1\% & 5 & 2 \\
\bottomrule
\end{tabular}
\end{table}

\paragraph{Scale does not change box selection.}
The two scales give the same lists. All 5,130 lists match between real and estimated scale in labels, number of boxes and
box scores (16.03 boxes per list on average in both). The estimated scene is the real scene times one factor per scene:
within a scene, the ratio of estimated to real coordinates is constant up to the 2-decimal rounding. That factor varies a lot
across scenes (percentiles 10/25/50/75/90 of est./real: 0.86 / 0.96 / 1.04 / 1.16 / 1.34). The \texttt{use\_scene} flags are
also identical between the real and estimated mix-v2-1 lists.

\paragraph{Counting.}
On all 565 counting questions, the named label has the same number of boxes at both scales. Only 18 answers change:
estimated scale is better on 9, real on 6, and 3 tie. For example, question 1987 has 6 boxes (ground truth 5) and gets 5 at
estimated and 4 at real scale; question 1990 has 9 boxes (ground truth 14) and gets 8 vs.\ 10. The model's answer depends
slightly on the coordinates in the prompt even though counting does not need them, and greedy decoding flips a few answers.
The 0.4-point difference on counting is noise, not an effect of scale. The other tasks that do not depend on scale (relative
distance, relative direction, appearance order, route planning) behave the same way: 0.3--4.1\% of answers change, with no
consistent direction.

\paragraph{Tasks that depend on scale.}
Absolute distance gains 2.2 points at real scale (59.7 vs.\ 57.5), as expected. Object size gains only 0.9 points,
although 57\% of its answers change. Room size does not change: under mix-v2-1 these questions are answered from the video
alone, so both scales score 72.8; with the list on every question, 70\% of the answers change but the score barely moves
(73.0 vs.\ 72.9).
; needs \usepackage{booktabs}
% and \usepackage{url}. Numbers from 3d-scene-lang/data/scenelang_vsibench_slam/scale_compare.py (gandalf).

\section{Real vs.\ estimated scale}
\label{sec:scale}

We compare the SAM~3 + VGGT-SLAM box lists at real scale (ground-truth scale from the camera fit) and at estimated scale
(from class size priors). We use the SFT Qwen3.5-4B Exp1 model
(\path{sr_vggt_novggt_qwen35_4b_reason190k_mixv2_1_objq/checkpoint-569}) on all 5,130 VSI-Bench questions. Each scale is run
with two inference rules: mix-v2-1 (video only for route planning, room size and questions with no named object) and
the list on every question (``full''). Table~\ref{tab:scale} gives the scores per task, together with the share of answers that
change between real and estimated scale (full rule) and which side is better on the changed answers.
Answers are extracted as the last option letter or number after \texttt{\textless/think\textgreater}. This makes the absolute values differ
slightly from the main tables (e.g.\ 59.7 vs.\ 60.8 on absolute distance for real mix-v2-1); the differences between runs
are what matter here.

\begin{table}[h]
\centering
\caption{Exp1 model with SAM~3 + VGGT-SLAM lists at real vs.\ estimated scale, per task (MRA or accuracy, \%).
Changed: answers that differ between real and estimated scale (full rule); est./real better: of the changed answers, how many
score higher at each scale.}
\label{tab:scale}
\small
\begin{tabular}{lrrrrrrrr}
\toprule
 & & \multicolumn{2}{c}{mix-v2-1} & \multicolumn{2}{c}{full} & \multicolumn{3}{c}{real $\rightarrow$ est.\ (full)} \\
\cmidrule(lr){3-4}\cmidrule(lr){5-6}\cmidrule(lr){7-9}
Task & $n$ & real & est. & real & est. & changed & est.\ better & real better \\
\midrule
\multicolumn{9}{l}{\emph{Depends on scale}} \\
Absolute distance & 834 & 59.7 & 57.5 & 59.7 & 57.5 & 52.3\% & 163 & 193 \\
Object size       & 953 & 76.4 & 75.5 & 76.4 & 75.5 & 56.9\% & 153 & 142 \\
Room size         & 288 & 72.8 & 72.8 & 72.9 & 73.0 & 70.5\% & 78 & 68 \\
\midrule
\multicolumn{9}{l}{\emph{Independent of scale}} \\
Counting                   & 565 & 75.4 & 75.8 & 75.4 & 75.8 & 3.2\% & 9 & 6 \\
Relative distance          & 710 & 69.2 & 68.7 & 69.2 & 68.7 & 1.4\% & 3 & 6 \\
Relative direction, easy   & 217 & 78.8 & 78.3 & 78.8 & 78.3 & 1.4\% & 1 & 2 \\
Relative direction, medium & 378 & 79.9 & 79.6 & 79.9 & 79.6 & 0.8\% & 1 & 2 \\
Relative direction, hard   & 373 & 72.9 & 72.9 & 72.9 & 72.9 & 1.3\% & 2 & 2 \\
Appearance order           & 618 & 80.4 & 80.4 & 80.7 & 80.7 & 0.3\% & 1 & 1 \\
Route planning             & 194 & 49.5 & 49.5 & 47.9 & 49.5 & 4.1\% & 5 & 2 \\
\bottomrule
\end{tabular}
\end{table}

\paragraph{Scale does not change box selection.}
The two scales give the same lists. All 5,130 lists match between real and estimated scale in labels, number of boxes and
box scores (16.03 boxes per list on average in both). The estimated scene is the real scene times one factor per scene:
within a scene, the ratio of estimated to real coordinates is constant up to the 2-decimal rounding. That factor varies a lot
across scenes (percentiles 10/25/50/75/90 of est./real: 0.86 / 0.96 / 1.04 / 1.16 / 1.34). The \texttt{use\_scene} flags are
also identical between the real and estimated mix-v2-1 lists.

\paragraph{Counting.}
On all 565 counting questions, the named label has the same number of boxes at both scales. Only 18 answers change:
estimated scale is better on 9, real on 6, and 3 tie. For example, question 1987 has 6 boxes (ground truth 5) and gets 5 at
estimated and 4 at real scale; question 1990 has 9 boxes (ground truth 14) and gets 8 vs.\ 10. The model's answer depends
slightly on the coordinates in the prompt even though counting does not need them, and greedy decoding flips a few answers.
The 0.4-point difference on counting is noise, not an effect of scale. The other tasks that do not depend on scale (relative
distance, relative direction, appearance order, route planning) behave the same way: 0.3--4.1\% of answers change, with no
consistent direction.

\paragraph{Tasks that depend on scale.}
Absolute distance gains 2.2 points at real scale (59.7 vs.\ 57.5), as expected. Object size gains only 0.9 points,
although 57\% of its answers change. Room size does not change: under mix-v2-1 these questions are answered from the video
alone, so both scales score 72.8; with the list on every question, 70\% of the answers change but the score barely moves
(73.0 vs.\ 72.9).
; needs \usepackage{booktabs}
% and \usepackage{url}. Numbers from 3d-scene-lang/data/scenelang_vsibench_slam/scale_compare.py (gandalf).

\section{Real vs.\ estimated scale}
\label{sec:scale}

We compare the SAM~3 + VGGT-SLAM box lists at real scale (ground-truth scale from the camera fit) and at estimated scale
(from class size priors). We use the SFT Qwen3.5-4B Exp1 model
(\path{sr_vggt_novggt_qwen35_4b_reason190k_mixv2_1_objq/checkpoint-569}) on all 5,130 VSI-Bench questions. Each scale is run
with two inference rules: mix-v2-1 (video only for route planning, room size and questions with no named object) and
the list on every question (``full''). Table~\ref{tab:scale} gives the scores per task, together with the share of answers that
change between real and estimated scale (full rule) and which side is better on the changed answers.
Answers are extracted as the last option letter or number after \texttt{\textless/think\textgreater}. This makes the absolute values differ
slightly from the main tables (e.g.\ 59.7 vs.\ 60.8 on absolute distance for real mix-v2-1); the differences between runs
are what matter here.

\begin{table}[h]
\centering
\caption{Exp1 model with SAM~3 + VGGT-SLAM lists at real vs.\ estimated scale, per task (MRA or accuracy, \%).
Changed: answers that differ between real and estimated scale (full rule); est./real better: of the changed answers, how many
score higher at each scale.}
\label{tab:scale}
\small
\begin{tabular}{lrrrrrrrr}
\toprule
 & & \multicolumn{2}{c}{mix-v2-1} & \multicolumn{2}{c}{full} & \multicolumn{3}{c}{real $\rightarrow$ est.\ (full)} \\
\cmidrule(lr){3-4}\cmidrule(lr){5-6}\cmidrule(lr){7-9}
Task & $n$ & real & est. & real & est. & changed & est.\ better & real better \\
\midrule
\multicolumn{9}{l}{\emph{Depends on scale}} \\
Absolute distance & 834 & 59.7 & 57.5 & 59.7 & 57.5 & 52.3\% & 163 & 193 \\
Object size       & 953 & 76.4 & 75.5 & 76.4 & 75.5 & 56.9\% & 153 & 142 \\
Room size         & 288 & 72.8 & 72.8 & 72.9 & 73.0 & 70.5\% & 78 & 68 \\
\midrule
\multicolumn{9}{l}{\emph{Independent of scale}} \\
Counting                   & 565 & 75.4 & 75.8 & 75.4 & 75.8 & 3.2\% & 9 & 6 \\
Relative distance          & 710 & 69.2 & 68.7 & 69.2 & 68.7 & 1.4\% & 3 & 6 \\
Relative direction, easy   & 217 & 78.8 & 78.3 & 78.8 & 78.3 & 1.4\% & 1 & 2 \\
Relative direction, medium & 378 & 79.9 & 79.6 & 79.9 & 79.6 & 0.8\% & 1 & 2 \\
Relative direction, hard   & 373 & 72.9 & 72.9 & 72.9 & 72.9 & 1.3\% & 2 & 2 \\
Appearance order           & 618 & 80.4 & 80.4 & 80.7 & 80.7 & 0.3\% & 1 & 1 \\
Route planning             & 194 & 49.5 & 49.5 & 47.9 & 49.5 & 4.1\% & 5 & 2 \\
\bottomrule
\end{tabular}
\end{table}

\paragraph{Scale does not change box selection.}
The two scales give the same lists. All 5,130 lists match between real and estimated scale in labels, number of boxes and
box scores (16.03 boxes per list on average in both). The estimated scene is the real scene times one factor per scene:
within a scene, the ratio of estimated to real coordinates is constant up to the 2-decimal rounding. That factor varies a lot
across scenes (percentiles 10/25/50/75/90 of est./real: 0.86 / 0.96 / 1.04 / 1.16 / 1.34). The \texttt{use\_scene} flags are
also identical between the real and estimated mix-v2-1 lists.

\paragraph{Counting.}
On all 565 counting questions, the named label has the same number of boxes at both scales. Only 18 answers change:
estimated scale is better on 9, real on 6, and 3 tie. For example, question 1987 has 6 boxes (ground truth 5) and gets 5 at
estimated and 4 at real scale; question 1990 has 9 boxes (ground truth 14) and gets 8 vs.\ 10. The model's answer depends
slightly on the coordinates in the prompt even though counting does not need them, and greedy decoding flips a few answers.
The 0.4-point difference on counting is noise, not an effect of scale. The other tasks that do not depend on scale (relative
distance, relative direction, appearance order, route planning) behave the same way: 0.3--4.1\% of answers change, with no
consistent direction.

\paragraph{Tasks that depend on scale.}
Absolute distance gains 2.2 points at real scale (59.7 vs.\ 57.5), as expected. Object size gains only 0.9 points,
although 57\% of its answers change. Room size does not change: under mix-v2-1 these questions are answered from the video
alone, so both scales score 72.8; with the list on every question, 70\% of the answers change but the score barely moves
(73.0 vs.\ 72.9).
; needs \usepackage{booktabs}
% and \usepackage{url}. Numbers from 3d-scene-lang/data/scenelang_vsibench_slam/scale_compare.py (gandalf).

\section{Real vs.\ estimated scale}
\label{sec:scale}

We compare the SAM~3 + VGGT-SLAM box lists at real scale (ground-truth scale from the camera fit) and at estimated scale
(from class size priors). We use the SFT Qwen3.5-4B Exp1 model
(\path{sr_vggt_novggt_qwen35_4b_reason190k_mixv2_1_objq/checkpoint-569}) on all 5,130 VSI-Bench questions. Each scale is run
with two inference rules: mix-v2-1 (video only for route planning, room size and questions with no named object) and
the list on every question (``full''). Table~\ref{tab:scale} gives the scores per task, together with the share of answers that
change between real and estimated scale (full rule) and which side is better on the changed answers.
Answers are extracted as the last option letter or number after \texttt{\textless/think\textgreater}. This makes the absolute values differ
slightly from the main tables (e.g.\ 59.7 vs.\ 60.8 on absolute distance for real mix-v2-1); the differences between runs
are what matter here.

\begin{table}[h]
\centering
\caption{Exp1 model with SAM~3 + VGGT-SLAM lists at real vs.\ estimated scale, per task (MRA or accuracy, \%).
Changed: answers that differ between real and estimated scale (full rule); est./real better: of the changed answers, how many
score higher at each scale.}
\label{tab:scale}
\small
\begin{tabular}{lrrrrrrrr}
\toprule
 & & \multicolumn{2}{c}{mix-v2-1} & \multicolumn{2}{c}{full} & \multicolumn{3}{c}{real $\rightarrow$ est.\ (full)} \\
\cmidrule(lr){3-4}\cmidrule(lr){5-6}\cmidrule(lr){7-9}
Task & $n$ & real & est. & real & est. & changed & est.\ better & real better \\
\midrule
\multicolumn{9}{l}{\emph{Depends on scale}} \\
Absolute distance & 834 & 59.7 & 57.5 & 59.7 & 57.5 & 52.3\% & 163 & 193 \\
Object size       & 953 & 76.4 & 75.5 & 76.4 & 75.5 & 56.9\% & 153 & 142 \\
Room size         & 288 & 72.8 & 72.8 & 72.9 & 73.0 & 70.5\% & 78 & 68 \\
\midrule
\multicolumn{9}{l}{\emph{Independent of scale}} \\
Counting                   & 565 & 75.4 & 75.8 & 75.4 & 75.8 & 3.2\% & 9 & 6 \\
Relative distance          & 710 & 69.2 & 68.7 & 69.2 & 68.7 & 1.4\% & 3 & 6 \\
Relative direction, easy   & 217 & 78.8 & 78.3 & 78.8 & 78.3 & 1.4\% & 1 & 2 \\
Relative direction, medium & 378 & 79.9 & 79.6 & 79.9 & 79.6 & 0.8\% & 1 & 2 \\
Relative direction, hard   & 373 & 72.9 & 72.9 & 72.9 & 72.9 & 1.3\% & 2 & 2 \\
Appearance order           & 618 & 80.4 & 80.4 & 80.7 & 80.7 & 0.3\% & 1 & 1 \\
Route planning             & 194 & 49.5 & 49.5 & 47.9 & 49.5 & 4.1\% & 5 & 2 \\
\bottomrule
\end{tabular}
\end{table}

\paragraph{Scale does not change box selection.}
The two scales give the same lists. All 5,130 lists match between real and estimated scale in labels, number of boxes and
box scores (16.03 boxes per list on average in both). The estimated scene is the real scene times one factor per scene:
within a scene, the ratio of estimated to real coordinates is constant up to the 2-decimal rounding. That factor varies a lot
across scenes (percentiles 10/25/50/75/90 of est./real: 0.86 / 0.96 / 1.04 / 1.16 / 1.34). The \texttt{use\_scene} flags are
also identical between the real and estimated mix-v2-1 lists.

\paragraph{Counting.}
On all 565 counting questions, the named label has the same number of boxes at both scales. Only 18 answers change:
estimated scale is better on 9, real on 6, and 3 tie. For example, question 1987 has 6 boxes (ground truth 5) and gets 5 at
estimated and 4 at real scale; question 1990 has 9 boxes (ground truth 14) and gets 8 vs.\ 10. The model's answer depends
slightly on the coordinates in the prompt even though counting does not need them, and greedy decoding flips a few answers.
The 0.4-point difference on counting is noise, not an effect of scale. The other tasks that do not depend on scale (relative
distance, relative direction, appearance order, route planning) behave the same way: 0.3--4.1\% of answers change, with no
consistent direction.

\paragraph{Tasks that depend on scale.}
Absolute distance gains 2.2 points at real scale (59.7 vs.\ 57.5), as expected. Object size gains only 0.9 points,
although 57\% of its answers change. Room size does not change: under mix-v2-1 these questions are answered from the video
alone, so both scales score 72.8; with the list on every question, 70\% of the answers change but the score barely moves
(73.0 vs.\ 72.9).
